
\begin{algorithm}[t]
\small
\caption{\textsc{UpdateSupport}(Hierarchy)}
\label{alg:hierarchy}
\KwIn{Remove r-clique se $R_{min}$, minimum heap $H_{\min}$, Clique Path Index $\tree$, $s$-support map $support$}
\KwOut{Updated $support$ and union-find structure}
\SetKwFunction{RemoveNode}{RemoveNode}
\SetKwFunction{CountPreEdges}{$s$-CliqueCount}
\SetKwProg{Fn}{Function}{:}{}
\SetKwFunction{UniteLevel}{UniteLevel}

\ForEach{clique path $\TPath$ containing $rc \in R_{min}$}{
    Update $support(c)$ and $H_{\min}$ for all $s$-clique $c \in \TPath$
    
    \colorbox{yellow!25}{\strut \UniteLevel{$rc,c,currentK$} for all $s$-clique $c \in \TPath$}
}

\Fn{\UniteLevel{$a, b, currentK$}}{
    \For{$k\gets currentK$ \KwTo $0$}{
        \If{not codeDisjointSets[k].unite(a, b)}{
            \Return;
        }
    }
}



\end{algorithm}
  
\section{Hierarchy Construction}
\label{sec:hierarchy}

% In addition to finding the nucleus decomposition of a graph, our proposed algorithms can also be used simply to find the hierarchy of the nuclei. The hierarchy is a tree structure where each node represents a nucleus, and the edges represent the containment relationship between nuclei. The root of the tree is the entire graph, and the leaves are the smallest nuclei.
Thanks to our framework, we can extend the nucleus decomposition to efficiently construct the hierarchy of nuclei. 
The hierarchy is represented as a tree structure, where each node corresponds to a nucleus and edges denote containment relationships between nuclei. 
The root of the tree represents the entire graph, while the leaves correspond to the smallest, most cohesive nuclei. 
% Hierarchy 的具体定义如下:
We formally define hierarchy construction in the context of nucleus decomposition.

\begin{definition}[$s$-Connectivity]\label{def:s-connect}
Given integers $0 < r < s$,  
two $r$-cliques $R$ and $R'$ in a graph (or subgraph) $H$ are said to be \emph{$s$-connected} 
if there exists a sequence of $r$-cliques
$
  R = R_1, R_2, \dots, R_t = R',
$
referred to as an \emph{$s$-clique path},  
such that for each $i \in \{1, \dots, t-1\}$ there exists an induced $s$-clique
$S_i \subseteq V(H)$ satisfying $R_i \cup R_{i+1} \subseteq S_i$.
\end{definition}


\begin{definition}[$s$-Connected Component]\label{def:s-component}
For integers $0 < r < s$,  
the \emph{$s$-connected component} of an $r$-clique $R$ is the induced subgraph
\[
  \mathrm{Comp}_s(R)
  = G\, \!\Bigl[\, \!\bigcup\nolimits_{R'\,:\,R'\text{ is }s\text{-connected to }R} R'\Bigr].
\]
\end{definition}

It is the maximal subgraph where all $r$-cliques are mutually $s$-connected, and such components partition all $r$-cliques in $G$.
% The $s$-connected components partition all $r$-cliques in $G$.

% \begin{definition}[$s$-Connected Component]\label{def:s-component}
% Given integers $0 < r < s$,  
% the \emph{$s$-connected component} of an $r$-clique $R$ is the induced subgraph
% \[
%   \mathrm{Comp}_s(R)
%   = G\!\Bigl[\,
%       \bigcup_{R' \text{ is } s\text{-connected to } R} R'
%     \Bigr].
% \]
% Equivalently, it is the maximal subgraph $H$ containing $R$ in which any two $r$-cliques are $s$-connected within $H$.  
% Since $s$-connectedness is an equivalence relation on the set of $r$-cliques, the $s$-connected components partition all $r$-cliques in $G$.
% \end{definition}



% 首先回顾一下Connected Component的定义. 给定两个r-clique $C_r$和$C_r'$，如果它们共享一个s-clique $C_s$，则称它们是s-connected的. 
% First, we revisit the definition of connected components. Given two $r$-cliques $C_r$ and $C_r'$, they are said to be $s$-connected if they share an $s$-clique $C_s$. An $r$-clique $C_r$ is said to be $s$-connected to another $r$-clique $C_r'$ if there exists a sequence of $r$-cliques $C_r=C_r^1, C_r^2, \ldots, C_r^k=C_r'$ such that each consecutive pair $C_r^i$ and $C_r^{i+1}$ are $s$-connected. The transitive closure of this relation partitions the set of all $r$-cliques into disjoint subsets called \emph{$(r,s)$-nucleus}. Each $(r,s)$-nucleus is a maximal set of $s$-connected $r$-cliques. 

Hierarchy can be constructed by identifying containment relationships among nuclei. Specifically, an $(r,s)$-nucleus $N$ is said to \emph{contain} another $(r',s')$-nucleus $N'$ if every $r'$-clique in $N'$ is also an $r$-clique in $N$, and $s' \ge s$. 
This containment relationship is represented as a directed edge from $N$ to $N'$ in the hierarchy tree. 
The \cpi can be used to efficiently determine whether two $r$-cliques are $s$-connected.

% by identifying the containment relationships among these nuclei. Specifically, an $(r,s)$-nucleus $N$ is said to contain another $(r',s')$-nucleus $N'$ if every $r'$-clique in $N'$ is also an $r$-clique in $N$, and $s' \geq s$. This containment relationship can be represented as a directed edge from $N$ to $N'$ in the hierarchy tree. We can use the \cpi to decide whether two $r$-cliques are $s$-connected:

\begin{theorem} [Clique Connectivity]
\label{thm:clique_connectivity}
Given Two $r$-cliques $C$ and $C'$ and a clique path $\TPath$ of the \cpi, if $C\subseteq \TPath$ and $C'\subseteq \TPath$, then $C$ and $C'$ are $s$-connected.
\end{theorem}

\begin{proof}
If $C\subseteq \TPath$ and $C'\subseteq \TPath$, then there exists a sequence of $r$-cliques $C=C^1, C^2, \ldots, C^k=C'$ such that each consecutive pair $C^i$ and $C^{i+1}$ are $s$-connected. This is because each $r$-clique in the sequence shares an $s$-clique with the next one, as they are all part of the same clique path $\TPath$. Therefore, by the definition of $s$-connectedness, $C$ and $C'$ are $s$-connected.
\end{proof} 

% 基于上述定理，我们可以使用多层并查集来识别所有的在不同$k$-核中的(r,s)-nucleus. 具体来说，我们使用一个多层并查集数据结构来维护每个$k$-核中的所有r-clique的连接组件, 并且在核心分解的过程中动态更新这些连接组件. 伪代码如算法\ref{alg:hierarchy}所示.

Based on the above theorem, we can use a multi-level disjoint set union (DSU) data structure to identify all $(r,s)$-nuclei. Specifically, we maintain a multi-level DSU where each level corresponds to a $k$-core, and within each level, we maintain the connected components of $r$-cliques. As we perform the core decomposition, we dynamically update these connected components. The pseudocode is shown in Algorithm~\ref{alg:hierarchy}.

\noindent
\stitle{Algorithm.}
% The procedure in Algorithm~\ref{alg:hierarchy} follows the base pipeline; we keep a stack of DSUs (one per level) to maintain $s$-components. For each batch $S_{\min}$ assigned level $k$, \UniteLevel applies unions on affected CPI paths and replicates each successful union to all $\ell\le k$; the downward short-circuit is justified by Lemma~\ref{lem:unite_shortcircuit}.
% 
The procedure in Algorithm~\ref{alg:hierarchy} follows the base pipeline; we keep a stack of DSUs (one per level) to maintain $s$-components. For each batch $S_{\min}$ assigned level $k$, \UniteLevel applies unions on affected \cpi paths and replicates each successful union to all levels at most $k$.

% \noindent
% \stitle{Correctness of \UniteLevel's short-circuit.}
% We maintain an array of disjoint-set unions (DSU) $\{\mathrm{DSU}[k]\}_{k\ge 0}$, one per $k$-level.
% The procedure \UniteLevel$(a,b,\mathrm{currentK})$ scans $k=\mathrm{currentK},\mathrm{currentK}-1,\ldots,0$ and calls $\mathrm{DSU}[k].\mathrm{unite}(a,b)$.
% It short-circuits (returns) at the first level $k^\*$ where $\mathrm{unite}$ returns \texttt{false}.
% We show that this early return is correct.
\stitle{Correctness.}
% \cite{firstNucleus}已经证明$k$-(r,s)-nucleus一定是$k-1$-(r,s)-nucleus的subgraph. 也就说明如果两个$r$-clique在$k$-level是connected的, 那么它们在$k-1$-level也一定是connected的.
A $k$-$(r,s)$-nucleus is always a subgraph of a $(k-1)$-$(r,s)$-nucleus~\cite{firstNucleus}. This implies that if two $r$-cliques are connected at level $k$, they must also remain connected at level $k-1$.
% 所以如果两个$s$-clique在剥离的过程中, 根据theorem \ref{thm:clique_connectivity}是连通的, 那么它们在所有更低的level也是连通的, 这也是Function \UniteLevel的行为. 
Therefore, if two $r$-cliques are found to be connected during peeling by Theorem~\ref{thm:clique_connectivity}, they remain connected at all lower levels. This property defines the behavior of the \UniteLevel function.
%
For the correctness of the early return (Line~6): 
if it were incorrect, there would exist a scenario in which two $r$-cliques are connected in the $k$-$(r,s)$-nucleus but not correctly connected in the $(k\!-\!1)$-$(r,s)$-nucleus. 
Such a situation is impossible because, during the peeling process, $c_{\min}$ is monotonically increasing, and each union operation is consistently propagated to all lower levels.
% 如果18行是错的, 那么则存在一种情况, 两个rclique在k-(r,s)-nuclear是连通的但是在(k-1)-(r,s)-nuclear并没有被正确的连通在一起. 这种情况不可能因为在剥离过程中, $c_{min}$是单调递增的, 而且每次union操作都会被持续地传递到更低的level上.
% The minimum clique support $c_{\min}$ is 
% monotonically non-decreasing
% throughout the peeling process. In the \textsc{UniteLevel} function, union operations are performed for all levels less than or equal to \texttt{currentK}.



% \begin{proof}[Proof]
% Initially every $\mathrm{DSU}[k]$ is a partition of singletons, so the invariant holds trivially.
% The only mutating operation is a successful union at some level $k$; by the procedure, the same union is propagated to all lower levels (i.e., levels $<k$) within the same call, and DSUs never split. Hence the invariant is preserved.
% If at level $k^{*}$ the call returns \texttt{false}, then $a$ and $b$ are already in the same set in $\mathrm{DSU}[k^{*}]$; by the invariant, they must also be in the same set for every level at most $k^{*}$.
% Therefore continuing the loop would only perform redundant (no-op) unions, and returning is correct.
% Conversely, when a union succeeds at some level $k$, continuing to lower levels is necessary to maintain the invariant.
% \end{proof}

% \begin{proof}[Proof]
% By construction, whenever a union succeeds at level $k$, the same union is immediately replicated to all lower levels within the same \UniteLevel call; DSUs never split, so this downward-replication invariant is preserved. If at some level $k^{*}$ the call returns \texttt{false}, then $a$ and $b$ are already in the same set in $\mathrm{DSU}[k^{*}]$, and by the invariant they are also in the same set at every lower level; further calls would be no-ops, hence returning is correct. When a union succeeds at a level, continuing downward is necessary precisely to maintain the invariant.
% \end{proof}

% \noindent\textit{Remark.}
% This mirrors the nesting property of $k$-levels in nucleus decomposition ($G_k\subseteq G_{k-1}$):
% once two $r$-cliques become connected at level $k$, they remain connected at all lower levels.
% Our downward replication enforces this monotonicity in the DSU array.

\stitle{Complexity.} The complexity is identical to that of nucleus decomposition.

